\chapter{Boolean Analysis — Gaussian Tail}
%%% generated by the blueprint pipeline from TCSlib.BooleanAnalysis.ThresholdFunctions.GaussianTail
\section{Overview}
This chapter develops the standard Gaussian density, cumulative distribution function, quantile, and total isoperimetric profile. It records Mills bounds and the asymptotic behavior of the profile near zero.
\section{Declarations}

\begin{definition}[Standard Gaussian probability measure]
\lean{BooleanAnalysis.ThresholdFunctions.standardGaussianMeasure}
\leanok
The standard Gaussian probability measure $\gamma$ on $\mathbb R$ is the measure with Lebesgue density $x\mapsto e^{-x^2/2}/\sqrt{2\pi}$. For every Borel set $A\subseteq\mathbb R$, it satisfies
\[
\gamma(A)=\int_A \frac{e^{-x^2/2}}{\sqrt{2\pi}}\,dx.
\]
\end{definition}

\begin{definition}[Standard Gaussian probability density]
\lean{BooleanAnalysis.ThresholdFunctions.standardGaussianPDF}
\leanok
The standard Gaussian density is the function $\varphi:\mathbb R\to\mathbb R$ defined, for every $x\in\mathbb R$, by
\[
\varphi(x)=\frac{e^{-x^2/2}}{\sqrt{2\pi}}.
\]
\end{definition}

\begin{definition}[Standard Gaussian distribution function]
\lean{BooleanAnalysis.ThresholdFunctions.standardGaussianCDF}
\leanok
\uses{BooleanAnalysis.ThresholdFunctions.standardGaussianMeasure}
Let $\gamma$ be the standard Gaussian probability measure on $\mathbb R$, with density $x\mapsto e^{-x^2/2}/\sqrt{2\pi}$. Its cumulative distribution function is the real-valued function defined for every $t\in\mathbb R$ by $\Phi(t)=\gamma(({-\infty},t])$.
\end{definition}

\begin{definition}[Complementary Gaussian distribution function]
\lean{BooleanAnalysis.ThresholdFunctions.standardGaussianTail}
\leanok
\uses{BooleanAnalysis.ThresholdFunctions.standardGaussianCDF}
Let $\Phi(t)=\int_{-\infty}^{t} e^{-x^2/2}/\sqrt{2\pi}\,dx$ be the standard Gaussian cumulative distribution function. Its complementary tail function is defined for every $t\in\mathbb R$ by $\overline\Phi(t)=1-\Phi(t)$.
\end{definition}

\begin{definition}[Total Gaussian quantile]
\lean{BooleanAnalysis.ThresholdFunctions.standardGaussianQuantile}
\leanok
\uses{BooleanAnalysis.ThresholdFunctions.standardGaussianCDF}
Let $\Phi(t)=\int_{-\infty}^{t}e^{-x^2/2}/\sqrt{2\pi}\,dx$ be the standard Gaussian cumulative distribution function. Define $Q:\mathbb R\to\mathbb R$ by
\[
Q(\alpha)=\inf\{t\in\mathbb R:\alpha\le\Phi(t)\},
\]
where the infimum is assigned the value $0$ when its set is empty or unbounded below. Thus $Q$ is the usual lower Gaussian quantile for $0<\alpha<1$, and $Q(\alpha)=0$ when $\alpha\le0$ or $\alpha\ge1$.
\end{definition}

\begin{definition}[Total Gaussian isoperimetric profile]
\lean{BooleanAnalysis.ThresholdFunctions.gaussianIsoperimetric}
\leanok
\uses{BooleanAnalysis.ThresholdFunctions.standardGaussianPDF,BooleanAnalysis.ThresholdFunctions.standardGaussianQuantile}
Let $\varphi(x)=e^{-x^2/2}/\sqrt{2\pi}$ and $\Phi(t)=\int_{-\infty}^{t}\varphi(x)\,dx$, and define $Q(\alpha)=\inf\{t\in\mathbb R:\alpha\le\Phi(t)\}$ for $0<\alpha<1$ and $Q(\alpha)=0$ otherwise. The Gaussian isoperimetric profile is the function $U:\mathbb R\to\mathbb R$ given by $U(\alpha)=\varphi(Q(\alpha))$. In particular its values at both endpoints and outside $[0,1]$ are $U(\alpha)=1/\sqrt{2\pi}$ for $\alpha\le0$ or $\alpha\ge1$.
\end{definition}

\begin{lemma}[Closed formula for the Gaussian density]
\lean{BooleanAnalysis.ThresholdFunctions.standardGaussianPDF_eq}
\leanok
\uses{BooleanAnalysis.ThresholdFunctions.standardGaussianPDF}
\difficulty{3}
Let $\varphi$ be the probability density of the standard Gaussian distribution $N(0,1)$. For every $x\in\mathbb R$, its value is
\[
\varphi(x)=\frac{1}{\sqrt{2\pi}}\exp(-x^2/2).
\]
\end{lemma}

\begin{lemma}[Strict positivity of the Gaussian density]
\lean{BooleanAnalysis.ThresholdFunctions.standardGaussianPDF_pos}
\leanok
\uses{BooleanAnalysis.ThresholdFunctions.standardGaussianPDF}
\difficulty{1}
Let $\varphi(x)=e^{-x^2/2}/\sqrt{2\pi}$ be the standard Gaussian density. For every $x\in\mathbb R$, one has $0<\varphi(x)$.
\end{lemma}

\begin{lemma}[Continuity of the Gaussian density]
\lean{BooleanAnalysis.ThresholdFunctions.continuous_standardGaussianPDF}
\leanok
\uses{BooleanAnalysis.ThresholdFunctions.standardGaussianPDF,BooleanAnalysis.ThresholdFunctions.standardGaussianPDF_eq}
\difficulty{3}
The standard Gaussian density $\varphi:\mathbb R\to\mathbb R$, defined by $\varphi(x)=e^{-x^2/2}/\sqrt{2\pi}$, is continuous on $\mathbb R$.
\end{lemma}

\begin{lemma}[Derivative of the Gaussian density]
\lean{BooleanAnalysis.ThresholdFunctions.hasDerivAt_standardGaussianPDF}
\leanok
\uses{BooleanAnalysis.ThresholdFunctions.standardGaussianPDF,BooleanAnalysis.ThresholdFunctions.standardGaussianPDF_eq}
\difficulty{4}
Let $\varphi(x)=e^{-x^2/2}/\sqrt{2\pi}$ be the standard Gaussian density. For every $x\in\mathbb R$, the function $\varphi$ is differentiable at $x$ and satisfies $\varphi'(x)=-x\varphi(x)$.
\end{lemma}

\begin{lemma}[Divergence of squares at negative infinity]
\lean{BooleanAnalysis.ThresholdFunctions.tendsto_sq_atBot}
\leanok
\difficulty{4}
For a real variable $x$ tending to $-\infty$, its square tends to $+\infty$:
\[
\lim_{x\to-\infty}x^2=+\infty.
\]
\end{lemma}

\begin{lemma}[Density decay at negative infinity]
\lean{BooleanAnalysis.ThresholdFunctions.tendsto_standardGaussianPDF_atBot}
\leanok
\uses{BooleanAnalysis.ThresholdFunctions.standardGaussianPDF,BooleanAnalysis.ThresholdFunctions.standardGaussianPDF_eq,BooleanAnalysis.ThresholdFunctions.tendsto_sq_atBot}
\difficulty{4}
Let $\varphi(x)=e^{-x^2/2}/\sqrt{2\pi}$ be the standard Gaussian density. Its limit at negative infinity is $\lim_{x\to-\infty}\varphi(x)=0$.
\end{lemma}

\begin{lemma}[Logarithm of the Gaussian density]
\lean{BooleanAnalysis.ThresholdFunctions.log_standardGaussianPDF}
\leanok
\uses{BooleanAnalysis.ThresholdFunctions.standardGaussianPDF,BooleanAnalysis.ThresholdFunctions.standardGaussianPDF_eq}
\difficulty{4}
Let $\varphi(x)=e^{-x^2/2}/\sqrt{2\pi}$ be the standard Gaussian density. For every $x\in\mathbb R$, its natural logarithm satisfies
\[
\log\varphi(x)=-\log\sqrt{2\pi}-\frac{x^2}{2}.
\]
\end{lemma}

\begin{lemma}[Integral formula for the Gaussian distribution]
\lean{BooleanAnalysis.ThresholdFunctions.standardGaussianCDF_eq_integral}
\leanok
\uses{BooleanAnalysis.ThresholdFunctions.standardGaussianCDF,BooleanAnalysis.ThresholdFunctions.standardGaussianMeasure,BooleanAnalysis.ThresholdFunctions.standardGaussianPDF}
\difficulty{3}
Let $\gamma$ be the standard Gaussian probability measure on $\mathbb R$, let $\Phi(t)=\gamma(({-\infty},t])$, and let $\varphi(x)=e^{-x^2/2}/\sqrt{2\pi}$. For every $t\in\mathbb R$, the cumulative probability equals the Lebesgue integral of the density:
\[
\Phi(t)=\int_{(-\infty,t]}\varphi(x)\,dx.
\]
\end{lemma}

\begin{lemma}[Gaussian distribution increments]
\lean{BooleanAnalysis.ThresholdFunctions.standardGaussianCDF_sub}
\leanok
\uses{BooleanAnalysis.ThresholdFunctions.standardGaussianCDF,BooleanAnalysis.ThresholdFunctions.standardGaussianCDF_eq_integral,BooleanAnalysis.ThresholdFunctions.standardGaussianPDF}
\difficulty{3}
Let $\varphi(x)=e^{-x^2/2}/\sqrt{2\pi}$ and $\Phi(t)=\int_{-\infty}^{t}\varphi(x)\,dx$. For any real $a,b$, with no ordering assumption, one has
\[
\Phi(b)-\Phi(a)=\int_a^b\varphi(x)\,dx,
\]
where the interval integral is signed, so that $\int_a^b=-\int_b^a$ when $b<a$.
\end{lemma}

\begin{lemma}[Derivative of the Gaussian distribution function]
\lean{BooleanAnalysis.ThresholdFunctions.hasDerivAt_standardGaussianCDF}
\leanok
\uses{BooleanAnalysis.ThresholdFunctions.continuous_standardGaussianPDF,BooleanAnalysis.ThresholdFunctions.standardGaussianCDF,BooleanAnalysis.ThresholdFunctions.standardGaussianCDF_sub,BooleanAnalysis.ThresholdFunctions.standardGaussianPDF}
\difficulty{4}
Let $\varphi(x)=e^{-x^2/2}/\sqrt{2\pi}$ and $\Phi(t)=\int_{-\infty}^{t}\varphi(x)\,dx$. For every $t\in\mathbb R$, the function $\Phi$ is differentiable at $t$ with derivative $\Phi'(t)=\varphi(t)$.
\end{lemma}

\begin{lemma}[Continuity of the Gaussian distribution function]
\lean{BooleanAnalysis.ThresholdFunctions.continuous_standardGaussianCDF}
\leanok
\uses{BooleanAnalysis.ThresholdFunctions.hasDerivAt_standardGaussianCDF,BooleanAnalysis.ThresholdFunctions.standardGaussianCDF}
\difficulty{1}
The standard Gaussian cumulative distribution function $\Phi:\mathbb R\to\mathbb R$, given by $\Phi(t)=\int_{-\infty}^{t}e^{-x^2/2}/\sqrt{2\pi}\,dx$, is continuous on $\mathbb R$.
\end{lemma}

\begin{lemma}[Strict increase of Gaussian cumulative probabilities]
\lean{BooleanAnalysis.ThresholdFunctions.standardGaussianCDF_strictMono}
\leanok
\uses{BooleanAnalysis.ThresholdFunctions.hasDerivAt_standardGaussianCDF,BooleanAnalysis.ThresholdFunctions.standardGaussianCDF,BooleanAnalysis.ThresholdFunctions.standardGaussianPDF_pos}
\difficulty{3}
Let $\Phi(t)=\int_{-\infty}^{t}e^{-x^2/2}/\sqrt{2\pi}\,dx$ be the standard Gaussian cumulative distribution function. It is strictly increasing: for every $a,b\in\mathbb R$ with $a<b$, one has $\Phi(a)<\Phi(b)$.
\end{lemma}

\begin{lemma}[Positive Gaussian cumulative probabilities]
\lean{BooleanAnalysis.ThresholdFunctions.standardGaussianCDF_pos}
\leanok
\uses{BooleanAnalysis.ThresholdFunctions.standardGaussianCDF,BooleanAnalysis.ThresholdFunctions.standardGaussianCDF_strictMono,BooleanAnalysis.ThresholdFunctions.standardGaussianMeasure}
\difficulty{3}
Let $\Phi(t)=\int_{-\infty}^{t}e^{-x^2/2}/\sqrt{2\pi}\,dx$ be the standard Gaussian cumulative distribution function. For every real $t$, one has $0<\Phi(t)$.
\end{lemma}

\begin{lemma}[Gaussian cumulative probabilities below one]
\lean{BooleanAnalysis.ThresholdFunctions.standardGaussianCDF_lt_one}
\leanok
\uses{BooleanAnalysis.ThresholdFunctions.standardGaussianCDF,BooleanAnalysis.ThresholdFunctions.standardGaussianCDF_strictMono,BooleanAnalysis.ThresholdFunctions.standardGaussianMeasure}
\difficulty{3}
Let $\Phi(t)=\int_{-\infty}^{t}e^{-x^2/2}/\sqrt{2\pi}\,dx$ be the standard Gaussian cumulative distribution function. For every real $t$, one has $\Phi(t)<1$.
\end{lemma}

\begin{lemma}[Attainment of interior Gaussian probabilities]
\lean{BooleanAnalysis.ThresholdFunctions.exists_standardGaussianCDF_eq}
\leanok
\uses{BooleanAnalysis.ThresholdFunctions.continuous_standardGaussianCDF,BooleanAnalysis.ThresholdFunctions.standardGaussianCDF,BooleanAnalysis.ThresholdFunctions.standardGaussianCDF_strictMono,BooleanAnalysis.ThresholdFunctions.standardGaussianMeasure}
\difficulty{5}
Let $\Phi(t)=\int_{-\infty}^{t}e^{-x^2/2}/\sqrt{2\pi}\,dx$ be the standard Gaussian cumulative distribution function. For every real $\alpha$ satisfying $0<\alpha<1$, there exists a real $t$ with $\Phi(t)=\alpha$.
\end{lemma}

\begin{lemma}[Quantile recovery from a cumulative value]
\lean{BooleanAnalysis.ThresholdFunctions.standardGaussianQuantile_eq}
\leanok
\uses{BooleanAnalysis.ThresholdFunctions.standardGaussianCDF,BooleanAnalysis.ThresholdFunctions.standardGaussianCDF_strictMono,BooleanAnalysis.ThresholdFunctions.standardGaussianQuantile}
\difficulty{4}
Let $\Phi(t)=\int_{-\infty}^{t}e^{-x^2/2}/\sqrt{2\pi}\,dx$, and let $Q(\alpha)$ be $\inf\{s\in\mathbb R:\alpha\le\Phi(s)\}$ for $0<\alpha<1$ and $0$ otherwise. For real $\alpha,t$ satisfying $\Phi(t)=\alpha$, one has $Q(\alpha)=t$.
\end{lemma}

\begin{lemma}[Cumulative probability of a Gaussian quantile]
\lean{BooleanAnalysis.ThresholdFunctions.standardGaussianCDF_quantile}
\leanok
\uses{BooleanAnalysis.ThresholdFunctions.exists_standardGaussianCDF_eq,BooleanAnalysis.ThresholdFunctions.standardGaussianCDF,BooleanAnalysis.ThresholdFunctions.standardGaussianQuantile,BooleanAnalysis.ThresholdFunctions.standardGaussianQuantile_eq}
\difficulty{4}
Let $\Phi(t)=\int_{-\infty}^{t}e^{-x^2/2}/\sqrt{2\pi}\,dx$, and let $Q(\alpha)$ be $\inf\{t\in\mathbb R:\alpha\le\Phi(t)\}$ for $0<\alpha<1$ and $0$ otherwise. For every $\alpha\in\mathbb R$ with $0<\alpha<1$, one has $\Phi(Q(\alpha))=\alpha$.
\end{lemma}

\begin{lemma}[Integrability of the Gaussian first moment]
\lean{BooleanAnalysis.ThresholdFunctions.integrable_id_mul_standardGaussianPDF}
\leanok
\uses{BooleanAnalysis.ThresholdFunctions.standardGaussianPDF}
\difficulty{4}
Let $\varphi(x)=e^{-x^2/2}/\sqrt{2\pi}$ be the standard Gaussian density. The function $x\mapsto x\varphi(x)$ is integrable over $\mathbb R$ with respect to Lebesgue measure.
\end{lemma}

\begin{lemma}[Truncated Gaussian first moment]
\lean{BooleanAnalysis.ThresholdFunctions.integral_Iic_id_mul_standardGaussianPDF}
\leanok
\uses{BooleanAnalysis.ThresholdFunctions.continuous_standardGaussianPDF,BooleanAnalysis.ThresholdFunctions.hasDerivAt_standardGaussianPDF,BooleanAnalysis.ThresholdFunctions.integrable_id_mul_standardGaussianPDF,BooleanAnalysis.ThresholdFunctions.standardGaussianPDF,BooleanAnalysis.ThresholdFunctions.tendsto_standardGaussianPDF_atBot}
\difficulty{4}
Let $\varphi(x)=e^{-x^2/2}/\sqrt{2\pi}$ be the standard Gaussian density. For every real $t$, its truncated first moment with respect to Lebesgue measure is
\[
\int_{(-\infty,t]}x\varphi(x)\,dx=-\varphi(t).
\]
\end{lemma}

\begin{lemma}[Upper Gaussian Mills bound]
\lean{BooleanAnalysis.ThresholdFunctions.standardGaussianCDF_le_mills}
\leanok
\uses{BooleanAnalysis.ThresholdFunctions.integrable_id_mul_standardGaussianPDF,BooleanAnalysis.ThresholdFunctions.integral_Iic_id_mul_standardGaussianPDF,BooleanAnalysis.ThresholdFunctions.standardGaussianCDF,BooleanAnalysis.ThresholdFunctions.standardGaussianCDF_eq_integral,BooleanAnalysis.ThresholdFunctions.standardGaussianPDF,BooleanAnalysis.ThresholdFunctions.standardGaussianPDF_pos}
\difficulty{5}
Let $\varphi(x)=e^{-x^2/2}/\sqrt{2\pi}$ and $\Phi(t)=\int_{-\infty}^{t}\varphi(x)\,dx$. For every real $t<0$, one has
\[
\Phi(t)\le\frac{\varphi(t)}{-t}.
\]
\end{lemma}

\begin{definition}[Mills lower-bound function]
\lean{BooleanAnalysis.ThresholdFunctions.millsF}
\leanok
\uses{BooleanAnalysis.ThresholdFunctions.standardGaussianPDF}
Let $\varphi(x)=e^{-x^2/2}/\sqrt{2\pi}$ be the standard Gaussian density. Define $F:\mathbb R\to\mathbb R$ by $F(0)=0$ and, for $x\ne0$, by
\[
F(x)=\varphi(x)(x^{-3}-x^{-1}).
\]
\end{definition}

\begin{definition}[Mills derivative integrand]
\lean{BooleanAnalysis.ThresholdFunctions.millsF'}
\leanok
\uses{BooleanAnalysis.ThresholdFunctions.standardGaussianPDF}
Let $\varphi(x)=e^{-x^2/2}/\sqrt{2\pi}$ be the standard Gaussian density. Define $H:\mathbb R\to\mathbb R$ by $H(0)=\varphi(0)$ and, for $x\ne0$, by
\[
H(x)=\varphi(x)(1-3x^{-4}).
\]
\end{definition}

\begin{lemma}[Derivative of the Mills lower-bound function]
\lean{BooleanAnalysis.ThresholdFunctions.hasDerivAt_millsF}
\leanok
\uses{BooleanAnalysis.ThresholdFunctions.hasDerivAt_standardGaussianPDF,BooleanAnalysis.ThresholdFunctions.millsF,BooleanAnalysis.ThresholdFunctions.millsF'}
\difficulty{4}
Let $\varphi(x)=e^{-x^2/2}/\sqrt{2\pi}$, let $F(0)=0$ and $F(x)=\varphi(x)(x^{-3}-x^{-1})$ for $x\ne0$, and let $H(0)=\varphi(0)$ and $H(x)=\varphi(x)(1-3x^{-4})$ for $x\ne0$. For every nonzero real $x$, the function $F$ is differentiable at $x$ with derivative $F'(x)=H(x)$.
\end{lemma}

\begin{lemma}[Measurability of the Mills derivative integrand]
\lean{BooleanAnalysis.ThresholdFunctions.measurable_millsF'}
\leanok
\uses{BooleanAnalysis.ThresholdFunctions.continuous_standardGaussianPDF,BooleanAnalysis.ThresholdFunctions.millsF'}
\difficulty{3}
Let $\varphi(x)=e^{-x^2/2}/\sqrt{2\pi}$, and define $H:\mathbb R\to\mathbb R$ by $H(0)=\varphi(0)$ and $H(x)=\varphi(x)(1-3x^{-4})$ for $x\ne0$. The function $H$ is Borel measurable.
\end{lemma}

\begin{lemma}[Integrability on negative lower half-lines]
\lean{BooleanAnalysis.ThresholdFunctions.integrableOn_millsF'}
\leanok
\uses{BooleanAnalysis.ThresholdFunctions.measurable_millsF',BooleanAnalysis.ThresholdFunctions.millsF',BooleanAnalysis.ThresholdFunctions.standardGaussianPDF,BooleanAnalysis.ThresholdFunctions.standardGaussianPDF_pos}
\difficulty{5}
Let $\varphi(x)=e^{-x^2/2}/\sqrt{2\pi}$, and define $H:\mathbb R\to\mathbb R$ by $H(0)=\varphi(0)$ and $H(x)=\varphi(x)(1-3x^{-4})$ for $x\ne0$. For every real $t<0$, the function $H$ is integrable on $(-\infty,t]$ with respect to Lebesgue measure.
\end{lemma}

\begin{lemma}[Cubic divergence at negative infinity]
\lean{BooleanAnalysis.ThresholdFunctions.tendsto_cube_atBot}
\leanok
\difficulty{4}
For a real variable $x$ tending to $-\infty$, its cube tends to $-\infty$:
\[
\lim_{x\to-\infty}x^3=-\infty.
\]
\end{lemma}

\begin{lemma}[Decay of the Mills lower-bound function]
\lean{BooleanAnalysis.ThresholdFunctions.tendsto_millsF_atBot}
\leanok
\uses{BooleanAnalysis.ThresholdFunctions.millsF,BooleanAnalysis.ThresholdFunctions.tendsto_cube_atBot,BooleanAnalysis.ThresholdFunctions.tendsto_standardGaussianPDF_atBot}
\difficulty{4}
Let $\varphi(x)=e^{-x^2/2}/\sqrt{2\pi}$, and define $F:\mathbb R\to\mathbb R$ by $F(0)=0$ and $F(x)=\varphi(x)(x^{-3}-x^{-1})$ for $x\ne0$. Its limit at negative infinity is $\lim_{x\to-\infty}F(x)=0$.
\end{lemma}

\begin{lemma}[Continuity of the Mills lower-bound function]
\lean{BooleanAnalysis.ThresholdFunctions.continuousAt_millsF}
\leanok
\uses{BooleanAnalysis.ThresholdFunctions.continuous_standardGaussianPDF,BooleanAnalysis.ThresholdFunctions.millsF}
\difficulty{4}
Let $\varphi(x)=e^{-x^2/2}/\sqrt{2\pi}$, and define $F:\mathbb R\to\mathbb R$ by $F(0)=0$ and $F(x)=\varphi(x)(x^{-3}-x^{-1})$ for $x\ne0$. For every nonzero real $t$, the function $F$ is continuous at $t$.
\end{lemma}

\begin{lemma}[Lower Gaussian Mills bound]
\lean{BooleanAnalysis.ThresholdFunctions.mills_le_standardGaussianCDF}
\leanok
\uses{BooleanAnalysis.ThresholdFunctions.continuousAt_millsF,BooleanAnalysis.ThresholdFunctions.hasDerivAt_millsF,BooleanAnalysis.ThresholdFunctions.integrableOn_millsF',BooleanAnalysis.ThresholdFunctions.millsF,BooleanAnalysis.ThresholdFunctions.millsF',BooleanAnalysis.ThresholdFunctions.standardGaussianCDF,BooleanAnalysis.ThresholdFunctions.standardGaussianCDF_eq_integral,BooleanAnalysis.ThresholdFunctions.standardGaussianPDF,BooleanAnalysis.ThresholdFunctions.standardGaussianPDF_pos,BooleanAnalysis.ThresholdFunctions.tendsto_millsF_atBot}
\difficulty{5}
Let $\varphi(x)=e^{-x^2/2}/\sqrt{2\pi}$ and $\Phi(t)=\int_{-\infty}^{t}\varphi(x)\,dx$, and define $F(0)=0$ and $F(x)=\varphi(x)(x^{-3}-x^{-1})$ for $x\ne0$. For every real $t<0$, one has
\[
F(t)\le\Phi(t).
\]
\end{lemma}

\begin{definition}[Signed Gaussian Mills ratio]
\lean{BooleanAnalysis.ThresholdFunctions.millsRatio}
\leanok
\uses{BooleanAnalysis.ThresholdFunctions.standardGaussianCDF,BooleanAnalysis.ThresholdFunctions.standardGaussianPDF}
Let $\varphi(x)=e^{-x^2/2}/\sqrt{2\pi}$ and $\Phi(t)=\int_{-\infty}^{t}\varphi(x)\,dx$. Define the real-valued function $R$ for every $t\in\mathbb R$ by
\[
R(t)=\frac{\Phi(t)(-t)}{\varphi(t)}.
\]
On $t<0$, this is the lower-tail Mills ratio $\Phi(t)|t|/\varphi(t)$.
\end{definition}

\begin{lemma}[Mills ratio upper bound]
\lean{BooleanAnalysis.ThresholdFunctions.millsRatio_le_one}
\leanok
\uses{BooleanAnalysis.ThresholdFunctions.millsRatio,BooleanAnalysis.ThresholdFunctions.standardGaussianCDF_le_mills,BooleanAnalysis.ThresholdFunctions.standardGaussianPDF_pos}
\difficulty{3}
Let $\varphi(x)=e^{-x^2/2}/\sqrt{2\pi}$, $\Phi(t)=\int_{-\infty}^{t}\varphi(x)\,dx$, and $R(t)=\Phi(t)(-t)/\varphi(t)$ for $t\in\mathbb R$. For every real $t<0$, one has $R(t)\le1$.
\end{lemma}

\begin{lemma}[Mills ratio lower bound]
\lean{BooleanAnalysis.ThresholdFunctions.one_sub_le_millsRatio}
\leanok
\uses{BooleanAnalysis.ThresholdFunctions.millsRatio,BooleanAnalysis.ThresholdFunctions.mills_le_standardGaussianCDF,BooleanAnalysis.ThresholdFunctions.standardGaussianCDF,BooleanAnalysis.ThresholdFunctions.standardGaussianPDF,BooleanAnalysis.ThresholdFunctions.standardGaussianPDF_pos}
\difficulty{4}
Let $\varphi(x)=e^{-x^2/2}/\sqrt{2\pi}$, $\Phi(t)=\int_{-\infty}^{t}\varphi(x)\,dx$, and $R(t)=\Phi(t)(-t)/\varphi(t)$ for $t\in\mathbb R$. For every real $t<0$, one has $1-t^{-2}\le R(t)$.
\end{lemma}

\begin{lemma}[Reciprocal-square decay at negative infinity]
\lean{BooleanAnalysis.ThresholdFunctions.tendsto_inv_sq_atBot}
\leanok
\uses{BooleanAnalysis.ThresholdFunctions.tendsto_sq_atBot}
\difficulty{1}
For a real variable $t$ tending to $-\infty$, the reciprocal of its square tends to zero:
\[
\lim_{t\to-\infty}\frac{1}{t^2}=0.
\]
\end{lemma}

\begin{lemma}[Gaussian Mills ratio limit]
\lean{BooleanAnalysis.ThresholdFunctions.tendsto_millsRatio}
\leanok
\uses{BooleanAnalysis.ThresholdFunctions.millsRatio,BooleanAnalysis.ThresholdFunctions.millsRatio_le_one,BooleanAnalysis.ThresholdFunctions.one_sub_le_millsRatio,BooleanAnalysis.ThresholdFunctions.tendsto_inv_sq_atBot}
\difficulty{4}
Let $\varphi(x)=e^{-x^2/2}/\sqrt{2\pi}$, $\Phi(t)=\int_{-\infty}^{t}\varphi(x)\,dx$, and $R(t)=\Phi(t)(-t)/\varphi(t)$ for $t\in\mathbb R$. The limit of $R$ at negative infinity is
\[
\lim_{t\to-\infty}R(t)=1.
\]
\end{lemma}

\begin{lemma}[Positive Mills ratio on negative arguments]
\lean{BooleanAnalysis.ThresholdFunctions.millsRatio_pos}
\leanok
\uses{BooleanAnalysis.ThresholdFunctions.millsRatio,BooleanAnalysis.ThresholdFunctions.standardGaussianCDF_pos,BooleanAnalysis.ThresholdFunctions.standardGaussianPDF_pos}
\difficulty{2}
Let $\varphi(x)=e^{-x^2/2}/\sqrt{2\pi}$, $\Phi(t)=\int_{-\infty}^{t}\varphi(x)\,dx$, and $R(t)=\Phi(t)(-t)/\varphi(t)$ for $t\in\mathbb R$. For every real $t<0$, one has $0<R(t)$.
\end{lemma}

\begin{definition}[Gaussian logarithmic tail scale]
\lean{BooleanAnalysis.ThresholdFunctions.gLog}
\leanok
\uses{BooleanAnalysis.ThresholdFunctions.standardGaussianCDF}
Let $\Phi(t)=\int_{-\infty}^{t}e^{-x^2/2}/\sqrt{2\pi}\,dx$ be the standard Gaussian cumulative distribution function. Define $L:\mathbb R\to\mathbb R$ for every real $t$ by
\[
L(t)=2\log\bigl(1/\Phi(t)\bigr),
\]
where $\log$ is the natural logarithm.
\end{definition}

\begin{lemma}[Positivity of the logarithmic tail scale]
\lean{BooleanAnalysis.ThresholdFunctions.gLog_pos}
\leanok
\uses{BooleanAnalysis.ThresholdFunctions.gLog,BooleanAnalysis.ThresholdFunctions.standardGaussianCDF,BooleanAnalysis.ThresholdFunctions.standardGaussianCDF_lt_one,BooleanAnalysis.ThresholdFunctions.standardGaussianCDF_pos}
\difficulty{4}
Let $\Phi(t)=\int_{-\infty}^{t}e^{-x^2/2}/\sqrt{2\pi}\,dx$ and $L(t)=2\log(1/\Phi(t))$, where $\log$ is the natural logarithm. For every real $t$, one has $0<L(t)$.
\end{lemma}

\begin{lemma}[Logarithmic expansion through the Mills ratio]
\lean{BooleanAnalysis.ThresholdFunctions.gLog_eq}
\leanok
\uses{BooleanAnalysis.ThresholdFunctions.gLog,BooleanAnalysis.ThresholdFunctions.log_standardGaussianPDF,BooleanAnalysis.ThresholdFunctions.millsRatio,BooleanAnalysis.ThresholdFunctions.standardGaussianCDF,BooleanAnalysis.ThresholdFunctions.standardGaussianCDF_pos,BooleanAnalysis.ThresholdFunctions.standardGaussianPDF,BooleanAnalysis.ThresholdFunctions.standardGaussianPDF_pos}
\difficulty{4}
Let $\varphi(x)=e^{-x^2/2}/\sqrt{2\pi}$, $\Phi(t)=\int_{-\infty}^{t}\varphi(x)\,dx$, $R(t)=\Phi(t)(-t)/\varphi(t)$, and $L(t)=2\log(1/\Phi(t))$. For every real $t<0$, with natural logarithms, one has
\[
L(t)=\bigl(-2\log R(t)+2\log(-t)+2\log\sqrt{2\pi}\bigr)+t^2.
\]
\end{lemma}

\begin{lemma}[Logarithm dominated by the square]
\lean{BooleanAnalysis.ThresholdFunctions.tendsto_log_neg_div_sq}
\leanok
\difficulty{4}
For a real variable $t$ tending to $-\infty$, with $\log$ denoting the natural logarithm, one has
\[
\lim_{t\to-\infty}\frac{\log(-t)}{t^2}=0.
\]
\end{lemma}

\begin{lemma}[Quadratic logarithmic tail asymptotics]
\lean{BooleanAnalysis.ThresholdFunctions.tendsto_gLog_div_sq}
\leanok
\uses{BooleanAnalysis.ThresholdFunctions.gLog,BooleanAnalysis.ThresholdFunctions.gLog_eq,BooleanAnalysis.ThresholdFunctions.millsRatio,BooleanAnalysis.ThresholdFunctions.tendsto_inv_sq_atBot,BooleanAnalysis.ThresholdFunctions.tendsto_log_neg_div_sq,BooleanAnalysis.ThresholdFunctions.tendsto_millsRatio}
\difficulty{5}
Let $\Phi(t)=\int_{-\infty}^{t}e^{-x^2/2}/\sqrt{2\pi}\,dx$ and $L(t)=2\log(1/\Phi(t))$, where $\log$ is the natural logarithm. Then
\[
\lim_{t\to-\infty}\frac{L(t)}{t^2}=1.
\]
\end{lemma}

\begin{lemma}[Square-root logarithmic tail asymptotics]
\lean{BooleanAnalysis.ThresholdFunctions.tendsto_sqrt_gLog_div}
\leanok
\uses{BooleanAnalysis.ThresholdFunctions.gLog,BooleanAnalysis.ThresholdFunctions.gLog_pos,BooleanAnalysis.ThresholdFunctions.tendsto_gLog_div_sq}
\difficulty{4}
Let $\Phi(t)=\int_{-\infty}^{t}e^{-x^2/2}/\sqrt{2\pi}\,dx$ and $L(t)=2\log(1/\Phi(t))$, where $\log$ is the natural logarithm. With the nonnegative square root, one has
\[
\lim_{t\to-\infty}\frac{\sqrt{L(t)}}{-t}=1.
\]
\end{lemma}

\begin{theorem}[Density relative to logarithmically scaled tail]
\lean{BooleanAnalysis.ThresholdFunctions.tendsto_standardGaussianPDF_div}
\leanok
\uses{BooleanAnalysis.ThresholdFunctions.gLog,BooleanAnalysis.ThresholdFunctions.gLog_pos,BooleanAnalysis.ThresholdFunctions.millsRatio,BooleanAnalysis.ThresholdFunctions.standardGaussianCDF,BooleanAnalysis.ThresholdFunctions.standardGaussianCDF_pos,BooleanAnalysis.ThresholdFunctions.standardGaussianPDF,BooleanAnalysis.ThresholdFunctions.standardGaussianPDF_pos,BooleanAnalysis.ThresholdFunctions.tendsto_millsRatio,BooleanAnalysis.ThresholdFunctions.tendsto_sqrt_gLog_div}
\difficulty{5}
Let $\varphi(x)=e^{-x^2/2}/\sqrt{2\pi}$, $\Phi(t)=\int_{-\infty}^{t}\varphi(x)\,dx$, and $L(t)=2\log(1/\Phi(t))$, with natural logarithm. Then
\[
\lim_{t\to-\infty}\frac{\varphi(t)}{\Phi(t)\sqrt{L(t)}}=1.
\]
\end{theorem}

\begin{theorem}[Quantile divergence at zero probability]
\lean{BooleanAnalysis.ThresholdFunctions.tendsto_standardGaussianQuantile}
\leanok
\uses{BooleanAnalysis.ThresholdFunctions.standardGaussianCDF,BooleanAnalysis.ThresholdFunctions.standardGaussianCDF_pos,BooleanAnalysis.ThresholdFunctions.standardGaussianCDF_quantile,BooleanAnalysis.ThresholdFunctions.standardGaussianCDF_strictMono,BooleanAnalysis.ThresholdFunctions.standardGaussianQuantile}
\difficulty{4}
Let $\Phi(t)=\int_{-\infty}^{t}e^{-x^2/2}/\sqrt{2\pi}\,dx$, and let $Q(\alpha)=\inf\{t\in\mathbb R:\alpha\le\Phi(t)\}$ for $0<\alpha<1$ and $Q(\alpha)=0$ otherwise. As a real $\alpha$ tends to zero through positive values, the quantile tends to negative infinity:
\[
\lim_{\alpha\to0^+}Q(\alpha)=-\infty.
\]
\end{theorem}

\begin{theorem}[Isoperimetric profile asymptotics near zero]
\lean{BooleanAnalysis.ThresholdFunctions.gaussianIsoperimetric_equiv_at_zero}
\leanok
\uses{BooleanAnalysis.ThresholdFunctions.gLog,BooleanAnalysis.ThresholdFunctions.gaussianIsoperimetric,BooleanAnalysis.ThresholdFunctions.standardGaussianCDF,BooleanAnalysis.ThresholdFunctions.standardGaussianCDF_quantile,BooleanAnalysis.ThresholdFunctions.standardGaussianQuantile,BooleanAnalysis.ThresholdFunctions.tendsto_standardGaussianPDF_div,BooleanAnalysis.ThresholdFunctions.tendsto_standardGaussianQuantile}
\difficulty{5}
Let $\varphi(x)=e^{-x^2/2}/\sqrt{2\pi}$, $\Phi(t)=\int_{-\infty}^{t}\varphi(x)\,dx$, and $Q(\alpha)=\inf\{t\in\mathbb R:\alpha\le\Phi(t)\}$ for $0<\alpha<1$, with $Q(\alpha)=0$ otherwise. Define $U(\alpha)=\varphi(Q(\alpha))$ for every real $\alpha$, so $U(\alpha)=1/\sqrt{2\pi}$ for $\alpha\le0$ or $\alpha\ge1$. As $\alpha\to0$ through positive values, this profile is asymptotically equivalent to $\alpha\sqrt{2\log(1/\alpha)}$; equivalently, with natural logarithm,
\[
\lim_{\alpha\to0^+}\frac{U(\alpha)}{\alpha\sqrt{2\log(1/\alpha)}}=1.
\]
\end{theorem}
